% =====================================================================
%  sections/05-tutorial.tex — AUTOTUTORIAL
%
%  Rozdział, w którym szablon opisuje sam siebie — na przykładach, które
%  sam demonstruje. Dołączany z main.tex.
%
%  PRZED ODDANIEM PRACY: usuń (lub zakomentuj) w main.tex linię
%      % =====================================================================
%  sections/05-tutorial.tex — AUTOTUTORIAL
%
%  Rozdział, w którym szablon opisuje sam siebie — na przykładach, które
%  sam demonstruje. Dołączany z main.tex.
%
%  PRZED ODDANIEM PRACY: usuń (lub zakomentuj) w main.tex linię
%      % =====================================================================
%  sections/05-tutorial.tex — AUTOTUTORIAL
%
%  Rozdział, w którym szablon opisuje sam siebie — na przykładach, które
%  sam demonstruje. Dołączany z main.tex.
%
%  PRZED ODDANIEM PRACY: usuń (lub zakomentuj) w main.tex linię
%      % =====================================================================
%  sections/05-tutorial.tex — AUTOTUTORIAL
%
%  Rozdział, w którym szablon opisuje sam siebie — na przykładach, które
%  sam demonstruje. Dołączany z main.tex.
%
%  PRZED ODDANIEM PRACY: usuń (lub zakomentuj) w main.tex linię
%      \include{sections/05-tutorial}
%  i zerknij na listę kontrolną z ostatniej podsekcji tego pliku.
% =====================================================================

\section*{Autotutorial, czyli szablon o sobie samym}
\addcontentsline{toc}{section}{Autotutorial, czyli szablon o sobie samym}
% ^ Właśnie obejrzałeś ręczny patent na rozdział bez numeru, ale z wpisem
%   w spisie treści. Tak samo zbudowany jest Wstęp, Podsumowanie i Dodatki
%   (a Numerowany zwykły rozdział dostajesz przez \section{Tytuł}).

Ten rozdział jest tu po to, żebyś TY\footnote{twój LLM} nie czytał dokumentacji pakietów
o drugiej w nocy. Każdy element, o którym mowa, powinien być gdzieś w szablonie
użyty na żywo — podglądaj PDF obok źródła, to najszybsza forma nauki.

\subsection*{Kompilacja}

Cały proces zamknięty jest w \verb|compile.sh|:

\begin{verbatim}
latexmk -bibtex -shell-escape -interaction=nonstopmode -pdfxe main.tex
\end{verbatim}

latexmk sam ustala, ile przebiegów potrzeba (spis treści i odsyłacze
\verb|\ref| „łapią" numery dopiero po drugim podejściu).
Dwie flagi zasługują na wyjaśnienie:
\begin{itemize}
	\item \verb|-shell-escape| — pakiet minted (kolorowanie kodu) odpala
	      zewnętrzny program \texttt{pygmentize}, a LaTeX domyślnie blokuje
	      uruchamianie obcych programów. Flaga jest zgodą. Jak pygmentize
	      nie masz: \verb|pip install Pygments|.
	\item \verb|-pdfxe| — silnik XeLaTeX, bo preambuła używa fontspec
	      i unicode-math (fonty systemowe, polskie znaki bez sztuczek).
\end{itemize}

Gdy kompilacja padnie: przeczytaj komunikat — pierwszy błąd się liczy,
kolejne to zwykle echa. Jeśli log nic sensownego nie mówi, usuń pliki
pomocnicze (\verb|latexmk -c|) i skompiluj od nowa.

\subsection*{Mapa projektu}

\begin{table}[H]
	\centering
	\caption{Mapa projektu. A tak wygląda tabela: środowisko
		\texttt{table} z podpisem \texttt{\textbackslash caption}
		i \texttt{\textbackslash label} zaraz po nim.}
	\label{tab:mapa-projektu}
	\begin{tabular}{|l|p{8.5cm}|}
		\hline
		\textbf{Plik}              & \textbf{Rola}                                  \\
		\hline
		\texttt{main.tex}          & szkielet: kolejność rozdziałów                 \\
		\texttt{core/packages.tex} & preambuła: pakiety, fonty, bibliografia        \\
		\texttt{core/macros.tex}   & dane pracy (temat, autor, promotor...) i makra \\
		\texttt{sections/*.tex}    & rozdziały — tu się pisze                       \\
		\texttt{img/}              & obrazki                                        \\
		\texttt{bibl.bib}          & baza źródeł do \texttt{\textbackslash cite}    \\
		\texttt{compile.sh}        & kompilacja jednym poleceniem                   \\
		\hline
	\end{tabular}
\end{table}

\subsection*{Dane pracy}

Temat (PL i EN), autor, promotor, data, typ pracy i angielskie Summary
to makra w \verb|core/macros.tex|. Rozkład strony tytułowej
(\verb|sections/0-head.tex|) odpowiada wymaganiom wydziału.

\subsection*{Pisanie rozdziałów}

\begin{itemize}
	\item \verb|\section{Tytuł}| — numerowany rozdział;
	\item rozdział bez numeru, ale w spisie treści (Wstęp, Podsumowanie,
	      Dodatki):
	      \verb|\section*{Wstęp}| plus \verb|\addcontentsline{toc}{section}{Wstęp}|;
	\item \verb|\subsection|, \verb|\subsubsection| — podział wewnątrz rozdziału.
\end{itemize}

Nowy rozdział to zwykle nowy plik w \verb|sections/| podpięty
\verb|\include| w \verb|main.tex| w odpowiednim miejscu listy.

\subsection*{Elementy składowe — gdzie zobaczyć na żywo}

Szablon zawiera działające przykłady wszystkiego, co typowej pracy
potrzebne:
\begin{itemize}
	\item rysunek kaczki z podpisem i źródłem: rozdział Wprowadzenie
	      (rysunek \ref{fig:duck} na str. \pageref{fig:duck};
	      tak, powtarzając 17 raz, to JEST KACZKA.
	\item tabela: badania eksperymentalne (tabela \ref{tab:xor-dataset}),
	\item wzór z odsyłaczem „(str. X)": implementacja, równanie
	      (\ref{eq:indeksliczenie}),
	\item listing z kolorowaniem składni (Python i Rust): analiza
	      istniejących rozwiązań oraz implementacja,
	\item algorytm w pseudokodzie i obrazek rysowany w TikZ-u: implementacja
	      (algorytm \ref{alg:tensor-iter}),
	\item cytowania \verb|\cite{...}|: analiza istniejących rozwiązań
	      i bibliografia na końcu dokumentu.
\end{itemize}

Żelazne zasady, które oszczędzają nerwów: \verb|\label| zawsze PO
\verb|\caption|, a każdy
rysunek i tabela dostają w podpisie źródło — „Źródło: opracowanie własne"
też jest źródłem.

\subsection*{Bibliografia}

Źródła dopisujesz do \verb|bibl.bib| i cytujesz \verb|\cite{klucz}| (wiele
naraz: \verb|\cite{a, b, c}|). Numery w bibliografii nadają się same, po
kolei cytowania — takie jest ustawienie \verb|sorting=none| w preambule.
Strony internetowe cytujemy z datą dostępu
(\verb|note = {[Dostęp dnia: 01.04.2025]}|). W \verb|bibl.bib| czeka po
jednym przykładzie każdego typu wpisu — od artykułu po stronę www.

\subsection*{Lista kontrolna przed oddaniem}

\begin{itemize}
	\item w \verb|main.tex| usunięta linia dołączająca ten tutorial,
	\item makra w \verb|core/macros.tex|: temat PL/EN, autor, promotor,
	      data, typ pracy, \textit{Summary} po angielsku,
	\item każdy rysunek i tabela mają podpis ze źródłem, ORAZ wspominasz
	      o nich w tekście (np. „na rysunku \ref{fig:duck} widać kaczkę. Naprawdę, to kaczka."),
	\item \verb|\label| po \verb|\caption| — odsyłacze mówią prawdę,
	\item dodatek o archiwum plików opisuje to, co faktycznie tam jest,
	\item końcowa kompilacja: log bez czerwonych linii, spis treści zgodny
	      z rzeczywistością,
	\item ortografia. Promotor czyta. Naprawdę.
\end{itemize}

Powodzenia.

%  i zerknij na listę kontrolną z ostatniej podsekcji tego pliku.
% =====================================================================

\section*{Autotutorial, czyli szablon o sobie samym}
\addcontentsline{toc}{section}{Autotutorial, czyli szablon o sobie samym}
% ^ Właśnie obejrzałeś ręczny patent na rozdział bez numeru, ale z wpisem
%   w spisie treści. Tak samo zbudowany jest Wstęp, Podsumowanie i Dodatki
%   (a Numerowany zwykły rozdział dostajesz przez \section{Tytuł}).

Ten rozdział jest tu po to, żebyś TY\footnote{twój LLM} nie czytał dokumentacji pakietów
o drugiej w nocy. Każdy element, o którym mowa, powinien być gdzieś w szablonie
użyty na żywo — podglądaj PDF obok źródła, to najszybsza forma nauki.

\subsection*{Kompilacja}

Cały proces zamknięty jest w \verb|compile.sh|:

\begin{verbatim}
latexmk -bibtex -shell-escape -interaction=nonstopmode -pdfxe main.tex
\end{verbatim}

latexmk sam ustala, ile przebiegów potrzeba (spis treści i odsyłacze
\verb|\ref| „łapią" numery dopiero po drugim podejściu).
Dwie flagi zasługują na wyjaśnienie:
\begin{itemize}
	\item \verb|-shell-escape| — pakiet minted (kolorowanie kodu) odpala
	      zewnętrzny program \texttt{pygmentize}, a LaTeX domyślnie blokuje
	      uruchamianie obcych programów. Flaga jest zgodą. Jak pygmentize
	      nie masz: \verb|pip install Pygments|.
	\item \verb|-pdfxe| — silnik XeLaTeX, bo preambuła używa fontspec
	      i unicode-math (fonty systemowe, polskie znaki bez sztuczek).
\end{itemize}

Gdy kompilacja padnie: przeczytaj komunikat — pierwszy błąd się liczy,
kolejne to zwykle echa. Jeśli log nic sensownego nie mówi, usuń pliki
pomocnicze (\verb|latexmk -c|) i skompiluj od nowa.

\subsection*{Mapa projektu}

\begin{table}[H]
	\centering
	\caption{Mapa projektu. A tak wygląda tabela: środowisko
		\texttt{table} z podpisem \texttt{\textbackslash caption}
		i \texttt{\textbackslash label} zaraz po nim.}
	\label{tab:mapa-projektu}
	\begin{tabular}{|l|p{8.5cm}|}
		\hline
		\textbf{Plik}              & \textbf{Rola}                                  \\
		\hline
		\texttt{main.tex}          & szkielet: kolejność rozdziałów                 \\
		\texttt{core/packages.tex} & preambuła: pakiety, fonty, bibliografia        \\
		\texttt{core/macros.tex}   & dane pracy (temat, autor, promotor...) i makra \\
		\texttt{sections/*.tex}    & rozdziały — tu się pisze                       \\
		\texttt{img/}              & obrazki                                        \\
		\texttt{bibl.bib}          & baza źródeł do \texttt{\textbackslash cite}    \\
		\texttt{compile.sh}        & kompilacja jednym poleceniem                   \\
		\hline
	\end{tabular}
\end{table}

\subsection*{Dane pracy}

Temat (PL i EN), autor, promotor, data, typ pracy i angielskie Summary
to makra w \verb|core/macros.tex|. Rozkład strony tytułowej
(\verb|sections/0-head.tex|) odpowiada wymaganiom wydziału.

\subsection*{Pisanie rozdziałów}

\begin{itemize}
	\item \verb|\section{Tytuł}| — numerowany rozdział;
	\item rozdział bez numeru, ale w spisie treści (Wstęp, Podsumowanie,
	      Dodatki):
	      \verb|\section*{Wstęp}| plus \verb|\addcontentsline{toc}{section}{Wstęp}|;
	\item \verb|\subsection|, \verb|\subsubsection| — podział wewnątrz rozdziału.
\end{itemize}

Nowy rozdział to zwykle nowy plik w \verb|sections/| podpięty
\verb|\include| w \verb|main.tex| w odpowiednim miejscu listy.

\subsection*{Elementy składowe — gdzie zobaczyć na żywo}

Szablon zawiera działające przykłady wszystkiego, co typowej pracy
potrzebne:
\begin{itemize}
	\item rysunek kaczki z podpisem i źródłem: rozdział Wprowadzenie
	      (rysunek \ref{fig:duck} na str. \pageref{fig:duck};
	      tak, powtarzając 17 raz, to JEST KACZKA.
	\item tabela: badania eksperymentalne (tabela \ref{tab:xor-dataset}),
	\item wzór z odsyłaczem „(str. X)": implementacja, równanie
	      (\ref{eq:indeksliczenie}),
	\item listing z kolorowaniem składni (Python i Rust): analiza
	      istniejących rozwiązań oraz implementacja,
	\item algorytm w pseudokodzie i obrazek rysowany w TikZ-u: implementacja
	      (algorytm \ref{alg:tensor-iter}),
	\item cytowania \verb|\cite{...}|: analiza istniejących rozwiązań
	      i bibliografia na końcu dokumentu.
\end{itemize}

Żelazne zasady, które oszczędzają nerwów: \verb|\label| zawsze PO
\verb|\caption|, a każdy
rysunek i tabela dostają w podpisie źródło — „Źródło: opracowanie własne"
też jest źródłem.

\subsection*{Bibliografia}

Źródła dopisujesz do \verb|bibl.bib| i cytujesz \verb|\cite{klucz}| (wiele
naraz: \verb|\cite{a, b, c}|). Numery w bibliografii nadają się same, po
kolei cytowania — takie jest ustawienie \verb|sorting=none| w preambule.
Strony internetowe cytujemy z datą dostępu
(\verb|note = {[Dostęp dnia: 01.04.2025]}|). W \verb|bibl.bib| czeka po
jednym przykładzie każdego typu wpisu — od artykułu po stronę www.

\subsection*{Lista kontrolna przed oddaniem}

\begin{itemize}
	\item w \verb|main.tex| usunięta linia dołączająca ten tutorial,
	\item makra w \verb|core/macros.tex|: temat PL/EN, autor, promotor,
	      data, typ pracy, \textit{Summary} po angielsku,
	\item każdy rysunek i tabela mają podpis ze źródłem, ORAZ wspominasz
	      o nich w tekście (np. „na rysunku \ref{fig:duck} widać kaczkę. Naprawdę, to kaczka."),
	\item \verb|\label| po \verb|\caption| — odsyłacze mówią prawdę,
	\item dodatek o archiwum plików opisuje to, co faktycznie tam jest,
	\item końcowa kompilacja: log bez czerwonych linii, spis treści zgodny
	      z rzeczywistością,
	\item ortografia. Promotor czyta. Naprawdę.
\end{itemize}

Powodzenia.

%  i zerknij na listę kontrolną z ostatniej podsekcji tego pliku.
% =====================================================================

\section*{Autotutorial, czyli szablon o sobie samym}
\addcontentsline{toc}{section}{Autotutorial, czyli szablon o sobie samym}
% ^ Właśnie obejrzałeś ręczny patent na rozdział bez numeru, ale z wpisem
%   w spisie treści. Tak samo zbudowany jest Wstęp, Podsumowanie i Dodatki
%   (a Numerowany zwykły rozdział dostajesz przez \section{Tytuł}).

Ten rozdział jest tu po to, żebyś TY\footnote{twój LLM} nie czytał dokumentacji pakietów
o drugiej w nocy. Każdy element, o którym mowa, powinien być gdzieś w szablonie
użyty na żywo — podglądaj PDF obok źródła, to najszybsza forma nauki.

\subsection*{Kompilacja}

Cały proces zamknięty jest w \verb|compile.sh|:

\begin{verbatim}
latexmk -bibtex -shell-escape -interaction=nonstopmode -pdfxe main.tex
\end{verbatim}

latexmk sam ustala, ile przebiegów potrzeba (spis treści i odsyłacze
\verb|\ref| „łapią" numery dopiero po drugim podejściu).
Dwie flagi zasługują na wyjaśnienie:
\begin{itemize}
	\item \verb|-shell-escape| — pakiet minted (kolorowanie kodu) odpala
	      zewnętrzny program \texttt{pygmentize}, a LaTeX domyślnie blokuje
	      uruchamianie obcych programów. Flaga jest zgodą. Jak pygmentize
	      nie masz: \verb|pip install Pygments|.
	\item \verb|-pdfxe| — silnik XeLaTeX, bo preambuła używa fontspec
	      i unicode-math (fonty systemowe, polskie znaki bez sztuczek).
\end{itemize}

Gdy kompilacja padnie: przeczytaj komunikat — pierwszy błąd się liczy,
kolejne to zwykle echa. Jeśli log nic sensownego nie mówi, usuń pliki
pomocnicze (\verb|latexmk -c|) i skompiluj od nowa.

\subsection*{Mapa projektu}

\begin{table}[H]
	\centering
	\caption{Mapa projektu. A tak wygląda tabela: środowisko
		\texttt{table} z podpisem \texttt{\textbackslash caption}
		i \texttt{\textbackslash label} zaraz po nim.}
	\label{tab:mapa-projektu}
	\begin{tabular}{|l|p{8.5cm}|}
		\hline
		\textbf{Plik}              & \textbf{Rola}                                  \\
		\hline
		\texttt{main.tex}          & szkielet: kolejność rozdziałów                 \\
		\texttt{core/packages.tex} & preambuła: pakiety, fonty, bibliografia        \\
		\texttt{core/macros.tex}   & dane pracy (temat, autor, promotor...) i makra \\
		\texttt{sections/*.tex}    & rozdziały — tu się pisze                       \\
		\texttt{img/}              & obrazki                                        \\
		\texttt{bibl.bib}          & baza źródeł do \texttt{\textbackslash cite}    \\
		\texttt{compile.sh}        & kompilacja jednym poleceniem                   \\
		\hline
	\end{tabular}
\end{table}

\subsection*{Dane pracy}

Temat (PL i EN), autor, promotor, data, typ pracy i angielskie Summary
to makra w \verb|core/macros.tex|. Rozkład strony tytułowej
(\verb|sections/0-head.tex|) odpowiada wymaganiom wydziału.

\subsection*{Pisanie rozdziałów}

\begin{itemize}
	\item \verb|\section{Tytuł}| — numerowany rozdział;
	\item rozdział bez numeru, ale w spisie treści (Wstęp, Podsumowanie,
	      Dodatki):
	      \verb|\section*{Wstęp}| plus \verb|\addcontentsline{toc}{section}{Wstęp}|;
	\item \verb|\subsection|, \verb|\subsubsection| — podział wewnątrz rozdziału.
\end{itemize}

Nowy rozdział to zwykle nowy plik w \verb|sections/| podpięty
\verb|\include| w \verb|main.tex| w odpowiednim miejscu listy.

\subsection*{Elementy składowe — gdzie zobaczyć na żywo}

Szablon zawiera działające przykłady wszystkiego, co typowej pracy
potrzebne:
\begin{itemize}
	\item rysunek kaczki z podpisem i źródłem: rozdział Wprowadzenie
	      (rysunek \ref{fig:duck} na str. \pageref{fig:duck};
	      tak, powtarzając 17 raz, to JEST KACZKA.
	\item tabela: badania eksperymentalne (tabela \ref{tab:xor-dataset}),
	\item wzór z odsyłaczem „(str. X)": implementacja, równanie
	      (\ref{eq:indeksliczenie}),
	\item listing z kolorowaniem składni (Python i Rust): analiza
	      istniejących rozwiązań oraz implementacja,
	\item algorytm w pseudokodzie i obrazek rysowany w TikZ-u: implementacja
	      (algorytm \ref{alg:tensor-iter}),
	\item cytowania \verb|\cite{...}|: analiza istniejących rozwiązań
	      i bibliografia na końcu dokumentu.
\end{itemize}

Żelazne zasady, które oszczędzają nerwów: \verb|\label| zawsze PO
\verb|\caption|, a każdy
rysunek i tabela dostają w podpisie źródło — „Źródło: opracowanie własne"
też jest źródłem.

\subsection*{Bibliografia}

Źródła dopisujesz do \verb|bibl.bib| i cytujesz \verb|\cite{klucz}| (wiele
naraz: \verb|\cite{a, b, c}|). Numery w bibliografii nadają się same, po
kolei cytowania — takie jest ustawienie \verb|sorting=none| w preambule.
Strony internetowe cytujemy z datą dostępu
(\verb|note = {[Dostęp dnia: 01.04.2025]}|). W \verb|bibl.bib| czeka po
jednym przykładzie każdego typu wpisu — od artykułu po stronę www.

\subsection*{Lista kontrolna przed oddaniem}

\begin{itemize}
	\item w \verb|main.tex| usunięta linia dołączająca ten tutorial,
	\item makra w \verb|core/macros.tex|: temat PL/EN, autor, promotor,
	      data, typ pracy, \textit{Summary} po angielsku,
	\item każdy rysunek i tabela mają podpis ze źródłem, ORAZ wspominasz
	      o nich w tekście (np. „na rysunku \ref{fig:duck} widać kaczkę. Naprawdę, to kaczka."),
	\item \verb|\label| po \verb|\caption| — odsyłacze mówią prawdę,
	\item dodatek o archiwum plików opisuje to, co faktycznie tam jest,
	\item końcowa kompilacja: log bez czerwonych linii, spis treści zgodny
	      z rzeczywistością,
	\item ortografia. Promotor czyta. Naprawdę.
\end{itemize}

Powodzenia.

%  i zerknij na listę kontrolną z ostatniej podsekcji tego pliku.
% =====================================================================

\section*{Autotutorial, czyli szablon o sobie samym}
\addcontentsline{toc}{section}{Autotutorial, czyli szablon o sobie samym}
% ^ Właśnie obejrzałeś ręczny patent na rozdział bez numeru, ale z wpisem
%   w spisie treści. Tak samo zbudowany jest Wstęp, Podsumowanie i Dodatki
%   (a Numerowany zwykły rozdział dostajesz przez \section{Tytuł}).

Ten rozdział jest tu po to, żebyś TY\footnote{twój LLM} nie czytał dokumentacji pakietów
o drugiej w nocy. Każdy element, o którym mowa, powinien być gdzieś w szablonie
użyty na żywo — podglądaj PDF obok źródła, to najszybsza forma nauki.

\subsection*{Kompilacja}

Cały proces zamknięty jest w \verb|compile.sh|:

\begin{verbatim}
latexmk -bibtex -shell-escape -interaction=nonstopmode -pdfxe main.tex
\end{verbatim}

latexmk sam ustala, ile przebiegów potrzeba (spis treści i odsyłacze
\verb|\ref| „łapią" numery dopiero po drugim podejściu).
Dwie flagi zasługują na wyjaśnienie:
\begin{itemize}
	\item \verb|-shell-escape| — pakiet minted (kolorowanie kodu) odpala
	      zewnętrzny program \texttt{pygmentize}, a LaTeX domyślnie blokuje
	      uruchamianie obcych programów. Flaga jest zgodą. Jak pygmentize
	      nie masz: \verb|pip install Pygments|.
	\item \verb|-pdfxe| — silnik XeLaTeX, bo preambuła używa fontspec
	      i unicode-math (fonty systemowe, polskie znaki bez sztuczek).
\end{itemize}

Gdy kompilacja padnie: przeczytaj komunikat — pierwszy błąd się liczy,
kolejne to zwykle echa. Jeśli log nic sensownego nie mówi, usuń pliki
pomocnicze (\verb|latexmk -c|) i skompiluj od nowa.

\subsection*{Mapa projektu}

\begin{table}[H]
	\centering
	\caption{Mapa projektu. A tak wygląda tabela: środowisko
		\texttt{table} z podpisem \texttt{\textbackslash caption}
		i \texttt{\textbackslash label} zaraz po nim.}
	\label{tab:mapa-projektu}
	\begin{tabular}{|l|p{8.5cm}|}
		\hline
		\textbf{Plik}              & \textbf{Rola}                                  \\
		\hline
		\texttt{main.tex}          & szkielet: kolejność rozdziałów                 \\
		\texttt{core/packages.tex} & preambuła: pakiety, fonty, bibliografia        \\
		\texttt{core/macros.tex}   & dane pracy (temat, autor, promotor...) i makra \\
		\texttt{sections/*.tex}    & rozdziały — tu się pisze                       \\
		\texttt{img/}              & obrazki                                        \\
		\texttt{bibl.bib}          & baza źródeł do \texttt{\textbackslash cite}    \\
		\texttt{compile.sh}        & kompilacja jednym poleceniem                   \\
		\hline
	\end{tabular}
\end{table}

\subsection*{Dane pracy}

Temat (PL i EN), autor, promotor, data, typ pracy i angielskie Summary
to makra w \verb|core/macros.tex|. Rozkład strony tytułowej
(\verb|sections/0-head.tex|) odpowiada wymaganiom wydziału.

\subsection*{Pisanie rozdziałów}

\begin{itemize}
	\item \verb|\section{Tytuł}| — numerowany rozdział;
	\item rozdział bez numeru, ale w spisie treści (Wstęp, Podsumowanie,
	      Dodatki):
	      \verb|\section*{Wstęp}| plus \verb|\addcontentsline{toc}{section}{Wstęp}|;
	\item \verb|\subsection|, \verb|\subsubsection| — podział wewnątrz rozdziału.
\end{itemize}

Nowy rozdział to zwykle nowy plik w \verb|sections/| podpięty
\verb|\include| w \verb|main.tex| w odpowiednim miejscu listy.

\subsection*{Elementy składowe — gdzie zobaczyć na żywo}

Szablon zawiera działające przykłady wszystkiego, co typowej pracy
potrzebne:
\begin{itemize}
	\item rysunek kaczki z podpisem i źródłem: rozdział Wprowadzenie
	      (rysunek \ref{fig:duck} na str. \pageref{fig:duck};
	      tak, powtarzając 17 raz, to JEST KACZKA.
	\item tabela: badania eksperymentalne (tabela \ref{tab:xor-dataset}),
	\item wzór z odsyłaczem „(str. X)": implementacja, równanie
	      (\ref{eq:indeksliczenie}),
	\item listing z kolorowaniem składni (Python i Rust): analiza
	      istniejących rozwiązań oraz implementacja,
	\item algorytm w pseudokodzie i obrazek rysowany w TikZ-u: implementacja
	      (algorytm \ref{alg:tensor-iter}),
	\item cytowania \verb|\cite{...}|: analiza istniejących rozwiązań
	      i bibliografia na końcu dokumentu.
\end{itemize}

Żelazne zasady, które oszczędzają nerwów: \verb|\label| zawsze PO
\verb|\caption|, a każdy
rysunek i tabela dostają w podpisie źródło — „Źródło: opracowanie własne"
też jest źródłem.

\subsection*{Bibliografia}

Źródła dopisujesz do \verb|bibl.bib| i cytujesz \verb|\cite{klucz}| (wiele
naraz: \verb|\cite{a, b, c}|). Numery w bibliografii nadają się same, po
kolei cytowania — takie jest ustawienie \verb|sorting=none| w preambule.
Strony internetowe cytujemy z datą dostępu
(\verb|note = {[Dostęp dnia: 01.04.2025]}|). W \verb|bibl.bib| czeka po
jednym przykładzie każdego typu wpisu — od artykułu po stronę www.

\subsection*{Lista kontrolna przed oddaniem}

\begin{itemize}
	\item w \verb|main.tex| usunięta linia dołączająca ten tutorial,
	\item makra w \verb|core/macros.tex|: temat PL/EN, autor, promotor,
	      data, typ pracy, \textit{Summary} po angielsku,
	\item każdy rysunek i tabela mają podpis ze źródłem, ORAZ wspominasz
	      o nich w tekście (np. „na rysunku \ref{fig:duck} widać kaczkę. Naprawdę, to kaczka."),
	\item \verb|\label| po \verb|\caption| — odsyłacze mówią prawdę,
	\item dodatek o archiwum plików opisuje to, co faktycznie tam jest,
	\item końcowa kompilacja: log bez czerwonych linii, spis treści zgodny
	      z rzeczywistością,
	\item ortografia. Promotor czyta. Naprawdę.
\end{itemize}

Powodzenia.
